\section{Abstraction}

Nếu decomposition giúp ta chia nhỏ vấn đề, còn pattern recognition giúp
ta tìm ra quy luật lặp, thì \textbf{abstraction} lại giúp ta loại bỏ các
chi tiết không thiết yếu để giữ lại cấu trúc bản chất của bài toán. Đây
là bước quan trọng trong phát triển phần mềm, vì thế giới thực luôn
rất phức tạp, còn máy tính thì chỉ làm việc hiệu quả khi bài toán được
biểu diễn dưới dạng các mô hình dữ liệu và các quan hệ rõ ràng~\cite{wing2006, spronck2023, openstax_ct}.

Trong BeVietnam, bài toán du lịch ngoài đời thực có thể chứa vô số chi
tiết khác nhau như cảm xúc cá nhân của du khách, mức độ nổi tiếng trên
mạng xã hội, độ đẹp của ảnh chụp hay rất nhiều yếu tố khó định lượng
khác. Nhưng để hệ thống có thể vận hành, ta buộc phải trừu tượng hóa các
thành phần đó thành các đối tượng và thuộc tính cốt lõi, tóm tắt trong
Định nghĩa~\ref{def:core-abstractions}.

\begin{mydef}{Bốn Lớp Trừu Tượng Cốt Lõi}{core-abstractions}
Hệ thống quy bài toán du lịch về bốn lớp trừu tượng:
\begin{enumerate}[label=\textbf{\arabic*.}, leftmargin=2.2cm, noitemsep]
    \item \textbf{\texttt{Place}} --- địa điểm với tọa độ, loại hình, mô tả
          và metadata phục vụ xếp hạng.
    \item \textbf{$\score$} --- mức độ phù hợp tại thời điểm hiện tại, một
          đại lượng tính được từ các yếu tố ngữ cảnh.
    \item \textbf{\texttt{Task}} --- một nút trong hành trình, có thứ tự,
          điều kiện mở khóa và tiêu chí hoàn thành.
    \item \textbf{\texttt{Capture}} --- bằng chứng thực địa, ánh xạ hành
          động ngoài đời của người dùng thành dữ liệu xác minh được.
\end{enumerate}
\end{mydef}

\textbf{Đầu tiên}, ta trừu tượng hóa khái niệm \textit{địa điểm du lịch}
thành thực thể \texttt{Place}. Một \texttt{Place} không còn là toàn bộ
trải nghiệm phong phú ngoài đời, mà được biểu diễn bằng các thuộc tính
như tên, tọa độ, loại hình, mô tả, hình ảnh, thông tin văn hóa và các
chỉ số phục vụ xếp hạng. Nhờ vậy máy tính mới có thể so sánh, lọc và
sắp xếp các địa điểm một cách tường minh. Trong quá trình làm đề tài,
lớp trừu tượng này còn giúp nhóm tách bạch ba tầng thông tin: dữ liệu
cốt lõi của POI, dữ liệu ngữ cảnh chỉ được nạp thêm khi tính score, và
nội dung diễn giải được sinh ra ở bước cuối để phục vụ người dùng.

\textbf{Kế đến}, ta trừu tượng hóa khái niệm \textit{mức độ phù hợp tại
thời điểm hiện tại} thành một \textbf{score} tổng hợp. Dù trên thực tế
khái niệm ``nên đi đâu bây giờ'' là một nhận định mang tính định tính,
nhưng để thuật toán có thể xử lý, ta phải biến nó thành một đại lượng có
thể tính toán từ các thành phần như khoảng cách, thời tiết, giao thông,
mật độ đông đúc và ưu tiên cá nhân. Trong tài liệu đặc tả, lớp trừu
tượng này đã được đẩy thêm một bước thành weighted score với các yếu tố
cultural value, distance, weather, traffic và estimated crowdedness cùng
tham gia vào quyết định cuối cùng~\cite{bevietnam_spec}. Đây là bước
trừu tượng hóa then chốt, vì nó buộc nhóm phải lượng hóa một câu hỏi rất
đời thường. Trước khi có lớp abstraction này, mọi trao đổi kiểu ``địa
điểm này có vẻ hợp hơn'' vẫn còn mang nặng trực giác. Khi đã có score,
nhóm mới có cơ sở để thảo luận về trọng số, về fallback và về
explainability theo đúng nghĩa kỹ thuật.

\textbf{Tiếp đó}, ta trừu tượng hóa \textit{hành trình khám phá văn hóa}
thành chuỗi các \texttt{Task} có thứ tự, điều kiện mở khóa và tiêu chí
hoàn thành. Điều này cực kỳ quan trọng, vì nếu không có lớp trừu tượng
này thì storyline sẽ chỉ còn là một đoạn văn giới thiệu dài chứ không
thể trở thành một cơ chế tương tác có thể điều khiển được bằng hệ thống.
Khi đưa vào thiết kế phần mềm, lớp trừu tượng này xuất hiện dưới dạng
kho nhiệm vụ (question pool), lộ trình Mission Path và các nút nhiệm vụ
có trạng thái mở khóa trên cả web lẫn mobile. Từ góc nhìn phân tích,
đây là bước biến một hành trình mang màu sắc kể chuyện thành một cấu
trúc có thể lưu trạng thái, xác minh tiến độ và mở khóa theo điều kiện
rõ ràng.

\textbf{Cuối cùng}, ta trừu tượng hóa hành động ngoài đời thật của người
dùng thành các \texttt{CaptureMetadata} và trạng thái xác minh. Nói cách
khác, việc ``đã thực sự tới nơi và chụp đúng đối tượng yêu cầu hay
chưa'' phải được ánh xạ thành dữ liệu đầu vào cho hệ thống kiểm tra. Từ
đó mà các agent hoặc service ở phía sau mới có thể đánh giá và phản hồi
một cách nhất quán. Nếu không làm bước này, toàn bộ ý tưởng geocaching
văn hóa sẽ chỉ còn là tuyên bố ở mức ý tưởng, bởi hệ thống sẽ không có
cách nào để nối hành động thật của người dùng với state số của task.

\begin{note}
Điểm cốt lõi của abstraction trong BeVietnam là chuyển những khái niệm
rất giàu ngữ nghĩa như ``một nơi đáng đi'', ``một hành trình khám phá''
hay ``đã thực sự hoàn thành nhiệm vụ'' thành các thực thể, trạng thái và
đại lượng mà hệ thống có thể xử lý được. Nếu không làm tốt bước này thì
mọi phần còn lại, từ scoring đến storyline hay capture verification, đều
sẽ chỉ dừng ở mức ý tưởng.
\end{note}
